\section{符号 AI、赛博系统与神经网络}

\subsection{两种 AI 观}

课程中一个非常重要的讨论，是“符号系统”和“神经系统”两种 AI 路线的差别。前者强调由符号、规则和结构组成的显式系统；后者强调通过数据和参数学习隐式表示。

\begin{figure}[H]
\centering
\includegraphics[width=0.95\textwidth]{slides-images/slide-024.jpg}
\caption{两种 AI 视角：符号 AI 与 cybernetics / 控制论传统\protect\footnotemark}
\end{figure}
\footnotetext{Slides 并列展示了符号 AI 与 cybernetics 的历史脉络。}

\begin{figure}[H]
\centering
\includegraphics[width=0.95\textwidth]{slides-images/slide-025.jpg}
\caption{Frank Rosenblatt 与 Perceptron：神经网络传统并非今天才出现，而是有很长的历史\protect\footnotemark}
\end{figure}
\footnotetext{Slides 展示了早期感知机和生物启发式连接主义。}

\begin{importantbox}{符号系统与神经系统各自擅长什么}
\begin{itemize}
    \item 符号系统擅长显式结构、可组合性和可控推理。
    \item 神经系统擅长从数据中学习、鲁棒泛化和容纳复杂模式。
    \item 今天的关键问题不是简单二选一，而是如何让两种能力互补。
\end{itemize}
\end{importantbox}

\subsection{传统语义学的启发}

在传统逻辑语义学里，意义常常被表示为“由词义和结构合成出来”的结果。一个经典例子是：

\begin{figure}[H]
\centering
\includegraphics[width=0.95\textwidth]{slides-images/slide-035.jpg}
\caption{传统视角下的语义组合：从句法树逐步构造句子的意义\protect\footnotemark}
\end{figure}
\footnotetext{Slides 用 “The red apple is on the table” 说明从句法到语义的组合过程。}

这种观点的核心是假设：先知道词的意义，再通过组合规则得到整句意义。它是形式语义学和很多逻辑化 NLP 系统的基础。

\begin{figure}[H]
\centering
\includegraphics[width=0.95\textwidth]{slides-images/slide-043.jpg}
\caption{神经模型是否真的提供可组合的意义函数？这是语义学和神经 NLP 之间的核心争论之一\protect\footnotemark}
\end{figure}
\footnotetext{Slides 直接提出了“Do neural models provide suitable meaning composition functions?”}

\subsection{本章小结}

符号系统和神经系统不是互相消灭的关系，而是两套不同的表达和推理范式。最终更实际的问题是：我们能否用神经模型承载足够多的结构，同时又保留足够强的统计学习能力。

%% ========== 意义与语义 ==========

